%%%%%%%%%%%%%%%%%%%%%%%%%%%%%%%%%%%%%%%%%%%%%%%%%
%%% Week 2: Dataset, Training, Main Results %%%%%
%%%%%%%%%%%%%%%%%%%%%%%%%%%%%%%%%%%%%%%%%%%%%%%%%

\chapter{Dataset, Training Strategy, and Main Results}\label{chap:data}

% STATUS: Week 2 (completed)

\section{Dataset Construction}\label{sec:dataset}

Our dataset comprises \textbf{604,589 images} derived from 151,147 unique source images, each processed into four compression variants: original (uncompressed), WhatsApp simulation, Instagram simulation, and screenshot simulation. The sources include:

\begin{itemize}
    \item \textbf{CIFAKE-Real}: Real images from the CIFAKE benchmark (CIFAR-10 source).
    \item \textbf{COCO}: Common Objects in Context~\cite{lin2014coco}, natural photographs covering 80 object categories.
    \item \textbf{FFHQ}: Flickr-Faces-HQ~\cite{karras2019ffhq}, 70,000 high-resolution face images.
    \item \textbf{SFHQ}: Synthetic Faces HQ~\cite{david2022sfhq}, synthetically generated face images.
    \item \textbf{Deepfake collections}: AI-generated face images from multiple GAN and diffusion architectures.
\end{itemize}

The overall class distribution is 57.9\% real and 42.1\% AI-generated images.

\subsection{Compression Variants}\label{sec:variants}

To simulate real-world distribution channels, each source image is processed through three social-media compression pipelines in addition to the original:

\begin{enumerate}
    \item \textbf{Original}: No compression applied; serves as the clean baseline.
    \item \textbf{WhatsApp}: JPEG recompression simulating WhatsApp's upload pipeline (quality reduction, resolution capping).
    \item \textbf{Instagram}: JPEG recompression simulating Instagram's processing (moderate quality, aspect ratio adjustments).
    \item \textbf{Screenshot}: Simulated screen capture (re-rendering artifacts, potential resolution changes).
\end{enumerate}

This gives $151{,}147 \times 4 = 604{,}588$ samples (rounding accounts for the actual 604,589 total).

\subsection{Data Split Strategy}\label{sec:split}

To prevent content leakage, the split is performed at the \textbf{image-id level} rather than at the individual file level: if a source image is assigned to the training set, all four of its compression variants (original, WhatsApp, Instagram, screenshot) remain in the training set and none can appear in the validation or test partitions.

Concretely, the splitting procedure works as follows:
\begin{enumerate}
    \item Extract all \textbf{unique image IDs} from the dataset (151,147 IDs).
    \item Apply \texttt{sklearn.train\_test\_split} on the unique IDs (not on individual samples) with stratification by class label.
    \item Split into train/validation/test with a 70/15/15\% ratio using seed 42.
    \item For each partition, include \textbf{all four compression variants} of each assigned image ID.
\end{enumerate}

This results in the following partition sizes:

\begin{table}[htbp]
    \centering
    \caption{Dataset split sizes. All four compression variants of each source image are assigned to the same partition.}
    \label{tab:split}
    \begin{tabular}{@{}lrrr@{}}
    \toprule
    \textbf{Partition} & \textbf{Samples} & \textbf{Real} & \textbf{AI-generated} \\
    \midrule
    Train      & 423,209 & 245,058 (57.9\%) & 178,151 (42.1\%) \\
    Validation &  90,688 &  52,508 (57.9\%) &  38,180 (42.1\%) \\
    Test       &  90,692 &  52,511 (57.9\%) &  38,181 (42.1\%) \\
    \midrule
    \textbf{Total} & \textbf{604,589} & \textbf{350,077} & \textbf{254,512} \\
    \bottomrule
    \end{tabular}
\end{table}

\section{Training Strategy}\label{sec:training}

\subsection{Phase 1: Embedding Extraction}

Three pre-trained backbone networks are used to extract penultimate-layer embeddings from each image:

\begin{itemize}
    \item \textbf{ResNet-50}~\cite{he2016deep}: Produces a 2,048-dimensional embedding vector. Deep residual network with skip connections; strong at capturing local texture patterns.
    \item \textbf{EfficientNet-B0}~\cite{tan2019efficientnet}: Produces a 1,280-dimensional embedding vector. Compound-scaled network optimised for parameter efficiency.
    \item \textbf{ViT-B/16}~\cite{dosovitskiy2021image}: Produces a 768-dimensional embedding vector. Vision Transformer with self-attention; captures global image context.
\end{itemize}

Embeddings are concatenated into a single feature vector:
\[
\mathbf{f} = [\mathbf{e}_1 \,\|\, \mathbf{e}_2 \,\|\, \mathbf{e}_3] \in \mathbb{R}^{4096}
\]
where $\mathbf{e}_1 \in \mathbb{R}^{2048}$ (ResNet-50), $\mathbf{e}_2 \in \mathbb{R}^{1280}$ (EfficientNet-B0), and $\mathbf{e}_3 \in \mathbb{R}^{768}$ (ViT-B/16).

\subsection{Phase 2: Meta-Classifier Training}

Two meta-classifiers are trained on the 4,096-dimensional fused vectors:

\paragraph{XGBoost meta-classifier.}
1,000 boosting rounds, maximum tree depth 8, learning rate 0.05, $L_1$ (lasso) and $L_2$ (ridge) regularisation of leaf weights, GPU-accelerated training, early stopping with patience 50 rounds.

\paragraph{MLP meta-classifier.}
Four-layer fully-connected neural network: $4096 \to 1024 \to 512 \to 256 \to 1$. Regularisation techniques: dropout (random neuron deactivation), label smoothing (soft targets instead of 0/1), mixup (convex interpolation of training pairs), SWA (Stochastic Weight Averaging). Optimiser: AdamW (Adam with decoupled weight decay). Early stopping with patience 15 epochs.

\paragraph{Reproducibility.}
Both meta-classifiers are trained with three seeds (42, 123, 7). Results are reported with bootstrap 95\% confidence intervals ($n = 1{,}000$ resamples). Hardware: GeForce RTX 5080 (16\,GB VRAM), AMD Ryzen 7 7800X3D, 32\,GB RAM.

\section{Baselines}\label{sec:baselines}

All baselines use the \textbf{same embeddings} on the \textbf{same train/validation/test splits}. They differ only in classifier architecture. This ensures a fair comparison where the only variable is the classification method.

\paragraph{External baselines.}
\begin{itemize}
    \item \textbf{Wang et al.}~\cite{wang2020cnn}: ResNet-50 pre-trained on ProGAN with blur and JPEG augmentation. Pre-trained weights evaluated directly on our test set.
    \item \textbf{Ojha et al.}~\cite{ojha2023towards}: Frozen CLIP ViT-L/14 representations with a trained linear probe layer. Zero-shot evaluation on our test set.
\end{itemize}

\paragraph{Internal baselines.}
\begin{itemize}
    \item \textbf{B1 -- LogReg + ResNet-50}: Logistic regression classifier trained on ResNet-50 embeddings only. Classical machine learning baseline.
    \item \textbf{B2 -- XGBoost + single CNN}: XGBoost trained on either ResNet-50 or EfficientNet-B0 embeddings alone. Evaluates the benefit of individual CNN backbones.
    \item \textbf{B3 -- XGBoost + ViT-B/16}: XGBoost trained on ViT-B/16 embeddings alone. Evaluates Transformer-based representations.
\end{itemize}

\paragraph{Metrics.}
Accuracy, Precision, Recall, F1 (harmonic mean of Precision and Recall), ROC-AUC (Area Under the Receiver Operating Characteristic curve), ECE (Expected Calibration Error, $M = 15$ bins), bootstrap 95\% CI, McNemar's test (paired comparison of classifier error rates), and DeLong's test (comparison of two AUC values).

\section{Main Results}\label{sec:main_results}

\begin{table}[htbp]
\centering
\caption{In-domain detection on 90,692 test images. Best in \textbf{bold}. $\downarrow$ = lower is better.}
\label{tab:main}
\setlength{\tabcolsep}{3.5pt}
\footnotesize
\begin{tabular}{@{}lcccccc@{}}
\toprule
\textbf{Method} & \textbf{Acc} & \textbf{Prec} & \textbf{Rec} & \textbf{F1} & \textbf{AUC} & \textbf{ECE}$\,\downarrow$ \\
\multicolumn{7}{@{}l}{\textit{External baselines}} \\
\quad Wang et al.   & 61.6 & 62.1 & 7.5 & 13.4 & 68.6 & 0.364 \\
\quad Ojha et al.   & 70.3 & 68.5 & 46.5 & 55.4 & 78.7 & 0.206 \\
\midrule
\multicolumn{7}{@{}l}{\textit{B1: Classical}} \\
\quad LogReg + ResNet-50     & 87.0 & 87.8 & 80.2 & 83.9 & 93.9 & 0.020 \\
\midrule
\multicolumn{7}{@{}l}{\textit{B2: Single-backbone XGBoost (CNN)}} \\
\quad XGBoost + ResNet-50        & 90.5 & 94.0 & 83.3 & 88.3 & 96.4 & 0.011 \\
\quad XGBoost + EfficientNet-B0  & 90.1 & 94.1 & 82.3 & 87.8 & 96.1 & 0.011 \\
\midrule
\multicolumn{7}{@{}l}{\textit{B3: Single-backbone XGBoost (Transformer)}} \\
\quad XGBoost + ViT-B/16         & 89.4 & 93.3 & 81.4 & 86.9 & 95.6 & 0.011 \\
\midrule
\multicolumn{7}{@{}l}{\textit{ImageTrust: 3-backbone fusion (ours)}} \\
\quad XGBoost meta-clf           & 88.7 & 90.5 & 81.9 & 85.9 & 96.0 & \textbf{0.016} \\
\quad \textbf{MLP meta-clf}      & \textbf{89.1} & 88.4 & \textbf{85.2} & \textbf{86.8} & \textbf{96.3} & 0.039 \\
\bottomrule
\end{tabular}
\end{table}

\subsection{Analysis of Results}

All methods in Table~\ref{tab:main}, including the external baselines, are evaluated on the \textbf{same} 90,692-image test split using identical metrics and preprocessing.

\paragraph{External baselines.}
The Wang et al.\ baseline (ResNet-50 pre-trained on ProGAN with blur and JPEG augmentation) achieves only 68.6\% AUC and 13.4\% F1 on our test set. The main failure mode is a severe \textbf{recall collapse} (7.5\%): the detector classifies nearly all images as real, yielding high precision (62.1\%) but missing 92.5\% of AI-generated images. This confirms that a detector trained on a single GAN architecture does not generalise to multi-source, multi-generator data.

Ojha et al.\ (frozen CLIP ViT-L/14 with a trained linear layer) significantly outperforms Wang, obtaining 78.7\% AUC and 55.4\% F1. That zero-shot predictions on frozen CLIP representations generalise better than Wang's supervised ResNet-50 (+10.1\% AUC) suggests large-scale vision--language pretraining captures more transferable forensic cues. However, Ojha still falls short of all our internal baselines: our simplest baseline (B1: LogReg) exceeds Ojha by +15.2\% AUC and +28.5\,pp F1.

\paragraph{Internal baselines vs.\ fusion.}
Our MLP meta-classifier obtains the highest AUC (0.963) and F1 (86.8\%). Most of the performance gain over XGBoost can be attributed to improved recall (85.2\% vs.\ 81.9\%). McNemar's test indicates a significant difference between classifiers ($\chi^2 = 20.76$, $p < 0.001$); DeLong's test does not find a significant difference in AUC ($p = 1.0$).

Single-backbone baselines obtain higher absolute accuracy because they learn from embeddings of original-only images, while fusion models learn from all four compression variants and sacrifice some accuracy for robustness to unseen degradations.

\paragraph{Key insight.}
Fusion does not increase absolute AUC over the best single-backbone baseline (96.3 vs 96.4); instead, it provides \textbf{better compression robustness} ($<$0.3\% AUC degradation under social-media recompression) and \textbf{calibrated uncertainty}, which single-backbone models lack.
